\subsection{有理性的判定准则}

对于 K 的每个子域 L，设 \(\operatorname{End}_{\mathrm{L}}(\mathrm{K})\) 表示把 K 视为 L 上的左向量空间时的自同态环；若 L 包含 \(K'\) ，则 \(\operatorname{End}_{\mathrm{L}}(\mathrm{K})\) 是 \(\operatorname{End}_{\mathrm{K}'}(\mathrm{K})\) 的一个子环。对于 \(\operatorname{End}_{\mathrm{K}'}(\mathrm{K})\) 的每个子集 M，存在 K 的一个包含 \(K'\) 且使得 M 含于 \(\operatorname{End}_{\mathrm{L}}(\mathrm{K})\) 的最大子域 L，即由这样的 \(\xi \in K\) 组成的集合：\(\phi(\xi\eta) = \xi.\phi(\eta)\) 对一切 \(\eta \in K\) 与一切 \(\phi \in M\) 成立（立即可验证此集合是一个环；另一方面，把前述关系中的 \(\eta\) 换成 \(\xi^{-1}\eta\) ，即得 \(\phi(\xi^{-1}\eta) = \xi^{-1}.\phi(\eta)\) （当 \(\xi \neq 0\) ）。我们称这个域为 M 在 K 中的中心化子，并记作 \(\chi(\mathcal{M})\) 。现在设 V 是带 \(K'\) -结构 \(V'\) 的右向量 K-空间。对一切 \(\phi \in \operatorname{End}_{\mathrm{K}'}(\mathrm{K})\) ，存在且仅存在 Z-模 V 的一个自同态 \(\phi_{V}\) ，使得 \(\phi_{\mathbf{V}}(x'.\xi) = x'.\phi(\xi)\) 对 \(x' \in \mathbf{V}'\) 与 \(\xi \in \mathbf{K}\) 成立：因为在第 1 号中曾定义过一个 Z-同构 \(\lambda\) ，它定义在 \(\mathbf{V}' \otimes_{\mathbb{K}'} \mathbf{K}\) 上且映到 V 上，把 \(x' \otimes \xi\) 映为 \(x'\) . \(\xi\) ，而 \(\phi_{\mathbf{V}}\) 必然等于 \(\lambda \circ (1_{\mathbf{V}'} \otimes \phi) \circ \lambda^{-1}\) 。

定理 1. 设 \(\mathcal{M}\) 是 \(\operatorname{End}_{\mathbb{K}'}(\mathbf{K})\) 的一个子集，而 \(\mathbf{L} = \chi(\mathcal{M})\) 是 \(\mathbf{K}\) 中作为 \(\mathcal{M}\) 的中心化子的那个子域。

\begin{enumerate}
\def\labelenumi{(\roman{enumi})}
\item
  设 V 是带 K'-结构的右向量 K-空间。向量 \(x \in V\) 在 L 上有理，当且仅当 \(\phi_{\mathbf{v}}(x.\eta) = x.\phi(\eta)\) 对一切 \(\phi \in M\) 与一切 \(\eta \in K\) 成立。V 的向量子 K-空间 W 在 L 上有理，当且仅当 \(\phi_{\mathbf{v}}(\mathbf{W}) \subset \mathbf{W}\) 对一切 \(\phi \in M\) 成立。
\item
  设 \(V_{1}\) , \(V_{2}\) 是两个各带 \(K^{\prime}\) -结构的右向量 K-空间。K-线性映射 f 从 \(V_{1}\) 到 \(V_{2}\) 在 L 上有理，当且仅当 \(f(\phi_{\mathrm{V}_{1}}(x_{1})) = \phi_{\mathrm{V}_{2}}(f(x_{1}))\) 对一切 \(x_{1} \in V_{1}\) 与一切 \(\phi \in M\) 成立。
\end{enumerate}

我们首先对 \(x\) 证明断言 (i)。设 \(B\) 是 \(V\) 的在 \(K'\) 上有理的一组基，并记 \(x = \sum_{b \in B} b \cdot \xi_b\) ；那么，对于 \(\phi \in \mathcal{M}\) 和 \(\eta \in K\) ，

\[
\phi_ {V} (x. \eta) - x. \phi (\eta) = \sum_ {b \in B} b. (\phi (\xi_ {b} \eta) - \xi_ {b}. \phi (\eta))
\]

因此，关系

\[
" \text { for   all } \phi \in \mathcal {M} \text { and   all } \eta \in K, \phi_ {V} (x. \eta) = x. \phi (\eta)"
\]

和

\[
" \text { for   all } \phi \in \mathcal {M}, \text { all } b \in \mathrm{B} \text { and   all } \eta \in \mathrm{K}, \phi (\xi_ {b} \eta) = \xi_ {b}. \phi (\eta)"
\]

是等价的。这些关系中的第二个意味着对所有 \(b \in \mathbf{B}\) , \(\xi_b \in \chi(\mathcal{M})\) ，这证明了(i)的第一个断言。

接下来证明 (ii)。\(f\) 在 \(L\) 上有理，当且仅当对所有 \(x_1' \in V_1\) 有理于 \(K', f(x_1')\) 是 \(V_2\) 中一个在 \(L\) 上有理的向量；这将蕴含 \(f(x_1)\) 在 \(L\) 上有理，只要 \(x_1\) 是 \(V_1\) 中在 \(L\) 上有理的向量，这样的向量乃是以 \(L\) 中系数对 \(K'\) 上有理的向量作的线性组合。由论证的第一部分，上述条件等价于关系式

\[
f (x _ {1} ^ {\prime}) \cdot \phi (\eta) = \phi_ {\mathrm{v} _ {2}} (f (x _ {1} ^ {\prime}) \cdot \eta) \quad \text { for } \phi \in \mathcal {M} \text { and } \eta \in \mathrm{K}\tag{4}
\]

也可以写作

\[
f (\phi_ {\mathrm{v} _ {1}} (x _ {1} ^ {\prime}. \eta)) = \phi_ {\mathrm{v} _ {2}} (f (x _ {1} ^ {\prime}. \eta)) \quad \text { for } \phi \in \mathcal {M} \text { and } \eta \in \mathrm{K}.\tag{5}
\]

由于 \(V_{1}\) 的每个元素都是以 K 中系数对 \(V_{1}\) 中在 \(K'\) 上有理的元素作的线性组合，条件 (5) 等价于

\[
f (\phi_ {\mathrm{v} _ {1}} (x _ {1})) = \phi_ {\mathrm{v} _ {2}} (f (x _ {1}))
\]

对所有 \(x_{1} \in V_{1}\) 与所有 \(\phi \in \mathcal{M}\) 。

最后，为证明 (i) 中的第二个断言，我们使用第 6 号引理 1：W 是某个 K-线性映射 \(g \colon \mathbb{W}_1 \to \mathbb{W}_2\) 的图像，而 W 在 L 上有理当且仅当映射 \(g\) 在 L 上有理（第 3 号，命题 4）。由 (ii)，\(g\) 在 L 上有理，当且仅当 \(g(\phi_{\mathbb{W}_1}(x_1)) = \phi_{\mathbb{W}_2}(g(x_1))\) 对所有 \(x_1 \in \mathbb{W}_1\) 与所有 \(\phi \in \mathcal{M}\) 成立；由于 \(\phi_{\mathbb{V}} = \phi_{\mathbb{W}_1} \times \phi_{\mathbb{W}_2}\) ，上述条件意味着 \(g\) 的图像 W 在 \(\phi_{\mathbb{V}}\) （对所有 \(\phi \in \mathcal{M}\) ）之下稳定。

